% REVISED: canonical-target reframe 2026-06-04
\begin{frontmatter}

%
\title{Cheap Signals, Costly Proof: The Reach and Limits of Award-Layer Screening in Cartel Enforcement\texorpdfstring{\footnote{We thank seminar participants at INSPER for comments; all remaining errors are our own.}}{}}

\author[insper]{Darcio Genicolo-Martins\corref{cor1}}
\ead{darciogm1@insper.edu.br}
\author[insper]{Paulo Furquim de Azevedo}
\affiliation[insper]{organization={INSPER},
  city={S\~ao Paulo}, country={Brazil}}
\cortext[cor1]{Corresponding author.}

%
\begin{abstract}
Cartel enforcement must allocate costly proof-producing effort
before legal proof exists. Using S\~ao Paulo's BEC procurement
platform and CADE adjudications, we ask how far routine award
records can rank where bid-level forensics should begin. We study
a minimal award-layer construct---persistent zero-win
participation, implemented as a frequent-loser flag---and validate
it against a reproducibly constructed, adjudication-anchored
cobidder label that never uses the screen itself. The decomposition is deliberately deflationary:
the cheap signal's apparent reach mostly reflects procurement
opportunity and case concentration; within comparable opportunity
sets the residual ordering is marginal at best, strict prospective
ranking fails outside incumbent pools, and exposure concentrates
in a few adjudicated cases. The transferable contribution is the
audit framework---label construction, opportunity adjustment,
timing discipline, case-composition and cost-recall
accounting---that shows where cheap screens stop and what
bid-layer follow-up costs. Liability remains in the richer
evidentiary record.
\end{abstract}

\end{frontmatter}

\noindent\textbf{Keywords:} cartel screening, public procurement, award-layer data, enforcement triage, incomplete observability, bid rigging

\noindent\textbf{JEL No.:} D44, D73, H57, K21, K42, L41

\bigskip
